% ======================================================================
%  Préambule commun — mémoire M1 Statistique
%  (chargé par main.tex)
% ======================================================================

% --- Encodage & langue -----------------------------------------------
\usepackage[utf8]{inputenc}
\usepackage[T1]{fontenc}
\usepackage[french]{babel}
\usepackage{lmodern}

% --- Mise en page ----------------------------------------------------
\usepackage[a4paper,margin=2.5cm]{geometry}
\usepackage{setspace}
\onehalfspacing

% Les subsubsections servent de titres NON numérotés (via \subsubsection*)
% pour nommer chaque méthode sans alourdir avec des 3.1.1, 3.1.2...
% On garde donc la numérotation standard jusqu'aux subsections.
\usepackage{parskip}
\usepackage{microtype}

% --- Mathématiques ---------------------------------------------------
\usepackage{amsmath,amssymb,amsthm}
\usepackage{bm}

% --- Tableaux & figures ----------------------------------------------
\usepackage{booktabs}
\usepackage{array}
\usepackage{multirow}
\usepackage{graphicx}
\graphicspath{{figures/}{../bloc1/03_resultats/}{../bloc3/03_resultats/}}
\usepackage{float}
\usepackage{caption}
\captionsetup{font=small,labelfont=bf}
\usepackage{subcaption}
\usepackage{pgfplots}
\pgfplotsset{compat=1.18}

% --- Couleurs & encadrés ---------------------------------------------
\usepackage{xcolor}
\definecolor{bleuunistra}{RGB}{0,61,110}
\definecolor{grisclair}{RGB}{242,242,242}
\usepackage[most]{tcolorbox}

% Encadré « conditions d'application »
\newtcolorbox{conditions}[1][]{
  colback=grisclair, colframe=bleuunistra, boxrule=0.5pt,
  fonttitle=\bfseries, title=Conditions d'application, #1}

% Encadré « limites de validité »
\newtcolorbox{limites}[1][]{
  colback=white, colframe=black!55, boxrule=0.5pt,
  fonttitle=\bfseries\itshape, title=Limites de validité, #1}

% --- Code -------------------------------------------------------------
\usepackage{listings}
\lstset{
  inputencoding=utf8,
  extendedchars=true,
  literate=%
    {é}{{\'e}}1 {è}{{\`e}}1 {ê}{{\^e}}1 {ë}{{\"e}}1
    {à}{{\`a}}1 {â}{{\^a}}1 {î}{{\^i}}1 {ï}{{\"i}}1
    {ô}{{\^o}}1 {û}{{\^u}}1 {ù}{{\`u}}1 {ç}{{\c c}}1
    {É}{{\'E}}1 {È}{{\`E}}1 {À}{{\`A}}1 {²}{{\textsuperscript{2}}}1,
  basicstyle=\ttfamily\small,
  breaklines=true,
  frame=single,
  backgroundcolor=\color{grisclair},
  language=Python,
  keywordstyle=\color{bleuunistra}\bfseries,
  showstringspaces=false,
  captionpos=b
}

% --- Liens & biblio --------------------------------------------------
\usepackage[hidelinks]{hyperref}
\usepackage{csquotes}
\usepackage[
  backend=biber,
  style=alphabetic,
  sorting=nyt
]{biblatex}
\addbibresource{biblio.bib}

% --- En-têtes --------------------------------------------------------
\usepackage{fancyhdr}
\setlength{\headheight}{13.6pt}      % évite le warning "headheight too small"
\addtolength{\topmargin}{-1.6pt}
\pagestyle{fancy}
\fancyhf{}
\fancyhead[L]{\small\nouppercase{\leftmark}}
\fancyfoot[C]{\small\thepage}
\renewcommand{\headrulewidth}{0.4pt}

% Les pages de debut de chapitre utilisent le style « plain » : on lui donne
% la meme pagination en bas, pour que le numero soit toujours au meme endroit.
\fancypagestyle{plain}{%
  \fancyhf{}%
  \fancyfoot[C]{\small\thepage}%
  \renewcommand{\headrulewidth}{0pt}%
}

% --- Raccourcis maison -----------------------------------------------
\newcommand{\edgevar}{\textit{edge}}
\newcommand{\Hzero}{\ensuremath{H_0}}
\newcommand{\Hun}{\ensuremath{H_1}}
\DeclareMathOperator{\Var}{Var}

% lmodern n'a pas de fonte « gras étendu + petites capitales » (bx/sc),
% demandée par certains titres : on la redirige vers les petites capitales
% normales pour supprimer le warning « Font shape bx/sc undefined ».
\AtBeginDocument{\DeclareFontShape{T1}{lmr}{bx}{sc}{<->ssub*lmr/m/sc}{}}
